\documentclass[11pt,a4paper,twocolumn]{article}

% === Core Packages ===
\usepackage[margin=2cm]{geometry}
\usepackage[utf8]{inputenc}
\usepackage[T1]{fontenc}
\usepackage{times}
\usepackage{amsmath,amssymb}
\usepackage{graphicx}
\usepackage{booktabs}
\usepackage{float}
\usepackage{enumitem}
\usepackage{tikz}
\usepackage{pgfplots}
\pgfplotsset{compat=1.18}
\usetikzlibrary{arrows.meta,shapes.geometric,positioning,calc,decorations.markings,fit,backgrounds}

% === Citation (IEEE style) ===
\usepackage[numbers,sort&compress]{natbib}
\bibliographystyle{IEEEtranN}

% === Hyperref ===
\usepackage[hidelinks]{hyperref}
\usepackage{xurl}

% === Settings ===
\setlength{\parindent}{1em}
\setlength{\parskip}{0.3em}

% === Title ===
\title{\textbf{Децентрализирана репутационна мрежа:\\Рамка за естествена обществена организация}}
\author{Pavel Kudrna\\Prague, Czechia\\Чернова v1 --- март 2026}
\date{}

\begin{document}
\maketitle

% ============================================================
% ABSTRACT
% ============================================================
\begin{abstract}
Чувството за справедливост---убеждението, че неправдите биват поправяни и заслугата се възнаграждава---е най-силната сплотяваща сила на обществото. Когато централизираните институции на управление не могат да осигурят това чувство, защото цената за отстояване на дадено право надхвърля стойността на самото право, обществената сплотеност се руши. Настоящите институционални структури не успяват да разрешат значителен клас спорове по систематично същия начин, по който централното планиране не успява да разпределя ресурси: чрез информационна асиметрия, липсващи обратни връзки и откъснатост на вземащите решения от последиците на техните решения. Тази статия представя Репутационна социална мрежа (RSN): децентрализиран, неподкупен, устойчив на цензура комуникационен слой, изграден върху децентрализирани идентификатори (DID), който позволява на псевдонимни участници да публикуват, проверяват и потребяват репутационна информация за поведение в реалния свят. Основният механизъм почива на три проектантски принципа: (1)~асиметрична структура на разходите, при която четенето на репутационни данни е евтино, но записът изисква проверяемо усилие; (2)~недетерминистичен протокол за избор на проверители, основан на consistent hashing, който остойностява радикалните твърдения чрез нарастващи разходи за итерация; и (3)~пазарно управляван модел на репутационни авторитети, при който частните проверители залагат натрупаната си репутация върху точността на твърденията, които подкрепят. Получаващата се архитектура произвежда емергентна меритокрация без централна координация и позволява постепенен, обратим път на миграция от съществуващите държавно администрирани системи.
\end{abstract}

% ============================================================
% 1. INTRODUCTION
% ============================================================
\section{Въведение}

\subsection{Проблемът на централизираното доверие}

Съвременното управление почива на едно фундаментално допускане: че централизираните институции могат да посредничат за доверието между непознати в мащаб. Съдилищата решават спорове. Регистрите удостоверяват собственост. Регулаторните органи налагат стандарти. Тези институции са проектирани в епоха, когато информацията е била оскъдна, комуникацията бавна, а делегирането на власт на централен орган е било единственият осъществим механизъм за координация. Централизираното управление обаче се проваля по същите структурни причини като централното планиране: информационна асиметрия, липсващи обратни връзки и откъснатост на вземащите решения от последиците на техните решения.

Допускането се проваля например когато цената на институционалното посредничество надхвърли стойността на спора. Да вземем наемател, който открива, че наетият му апартамент няма законово изисквания сертификат за електрическа безопасност---състояние, което представлява непосредствена опасност за живота. Законовото право на наемателя да се откаже от договора е недвусмислено. Все пак цената за отстояване на това право чрез съдебната система надхвърля паричната стойност на депозита и дължимите обезщетения, дори включително потенциалната законова компенсация~\cite{czcourtfees}. Рационалният отговор е да се поеме загубата и да се излезе.

Това не е патологичен граничен случай. Това е стандартното преживяване за значителен клас спорове, при които вредата е реална, но паричните залози падат под прага, при който институционалното принудително изпълнение става икономически рационално. Резултатът е система, която структурно облагодетелства недобросъвестните: тези, които вредят на другите в обеми под този праг, могат да действат безнаказано, защото цената за търсене на справедливост пада върху ощетената страна.

Горният пример е взет от правната среда на автора---чешката континентална правна система. В други правни традиции---прецедентното common law, обичайното право, смесени системи---справедливостта се проваля по различни начини и в различна степен, в зависимост от културните и процедурните особености на юрисдикцията. Читателите вероятно ще разпознаят еквивалентни примери от собствения си контекст. Структурно универсален е моделът: спорове, чиято стойност пада под прага на икономически рационалното принудително изпълнение, завършват с това, че загубата бива поета от ощетената страна. RSN не обещава съвършена справедливост---тя просто намалява структурното предимство, на което се радват недобросъвестните, когато институционалното принудително изпълнение е неикономично.

\subsection{Защо съществуващите решения не са достатъчни}

\textbf{Рамките за децентрализирана идентичност (DID)}~\cite{did2022} предоставят криптографски механизми за самостоятелно управление на идентичността, но се фокусират преди всичко върху проверката на идентичността, без да разглеждат по-широкия въпрос за поведенческата репутация.

\textbf{Soulbound токени (SBT)}~\cite{weyl2022} улавят прозрението, че репутацията е непрехвърляема, но разчитат на блокчейн инфраструктура с ограничения по мащабируемост и не разглеждат икономическите стимули, управляващи въвеждането на информация. Свързани концепции като меритни токени (напр.\ Liberland) споделят подобна философия, но остават обвързани с конкретни юрисдикционни рамки.

\textbf{Децентрализирани автономни организации (DAO)}~\cite{buterin2014} реализират управление чрез смарт договори, но възпроизвеждат плутократични динамики, при които влиянието е пропорционално на капитала. Квадратичното гласуване~\cite{buterin2019} частично смекчава това, но ограничава практическото приемане. DAO се сблъскват и с фундаменталния \emph{oracle проблем}: смарт договорите не могат надеждно да проверяват състояния от реалния свят без доверен външен източник, а този проблем няма общо решение в рамките на блокчейн архитектурата.

Нито една съществуваща система не съчетава децентрализация, устойчивост на цензура, псевдонимност, проверяема цена на въвеждане и емергентен консенсус.

\subsection{Приноси}

Настоящата статия прави следните приноси:
\begin{enumerate}[nosep]
\item Дефиниция на \textbf{Репутационната социална мрежа (RSN)}, децентрализиран протокол, разширяващ рамката DID, за да поддържа двустранен обмен на репутация.
\item \textbf{Недетерминистичен протокол за избор на проверители}, основан на consistent hashing~\cite{karger1997}.
\item \textbf{Принципът на скъпия радикализъм}, който остойностява твърденията според тяхното разстояние от мрежовия консенсус чрез три взаимоусилващи се разходни канала.
\item Модел на \textbf{репутационен авторитет} с асиметрични стимули, при които фалшивото подкрепяне струва катастрофално повече от легитимните такси.
\item \textbf{Постепенен път на миграция} чрез паралелна фискална инфраструктура.
\end{enumerate}

% ============================================================
% 2. DESIGN PRINCIPLES
% ============================================================
\section{Проектантски принципи}

Архитектурата на RSN е ограничена от пет непреговаряеми проектантски принципа.

\textbf{Приемственост и еволюция.} Системата съществува съвместно със съществуващите институции по време на преходен период с неопределена продължителност. Преходът трябва да бъде обратим на всеки етап.

\textbf{Доброволност и отговорност.} Участието е доброволно, но доброволността е свързана с отговорност: участник, който поеме контрол над институционална функция, едновременно поема и съпътстващите я задължения.

\textbf{Многообразие и културна неутралност.} Консенсусът възниква от двустранни взаимодействия, а не от предварително зададени правила, наложени от проектантите на системата.

\textbf{Устойчивост и съпротива срещу цензура.} Цената за цензуриране на мрежата трябва да надхвърля цената за търпимостта към нея. Само пълен тоталитаризъм---на разорителна политическа и икономическа цена---може да потисне системата.

\textbf{Консенсус и принудително изпълнение.} Системата включва човешките склонности към дребнавост, злоба и възмездно удовлетворение като носещи функции. Простъпката не е забранена; тя е остойностена.

% ============================================================
% 3. SYSTEM OVERVIEW
% ============================================================
\section{Преглед на системата}

\subsection{Репутационна социална мрежа}

RSN е децентрализиран, peer-to-peer комуникационен слой, в който участниците обменят проверяема информация за поведение в реалния свят. Определящите ѝ свойства са: децентрализирана и нецензурируема инфраструктура; псевдонимно участие чрез DID; асиметрична структура на разходите (четенето е евтино, записът е скъп); криптографска проверяемост; и \textbf{поддръжка на стандарти}---машинно четими формати за информацията, разпространявана през мрежата, за услугите, предлагани от авторитетите (съставяне на твърдения, подкрепяне, свидетелска услуга), и за формата на политиката, включително правилата за оценка на репутацията на насрещната страна и за оценка на риска от несъответствие на политиката (между декларираното и действителното поведение).

\subsection{Архитектурни компоненти}

RSN се състои от четири основни компонента (Фиг.~\ref{fig:architecture}):
\begin{enumerate}[nosep]
\item \textbf{DID слой.} Идентичности на участниците с декларирани правила, репутационни метаданни и мрежово маршрутизиране.
\item \textbf{Протокол за твърдения.} Структуриран формат за публикуване на твърдения за събития от реалния свят.
\item \textbf{Мрежа за проверка.} Децентрализиран протокол за назначаване на проверители чрез consistent hashing.
\item \textbf{Слой за агрегиране на репутация.} Услуги, които произвеждат четими за човек репутационни обобщения.
\end{enumerate}

\begin{figure}[ht]
\centering
\begin{tikzpicture}[
    layer/.style={draw, minimum height=0.7cm, align=center, font=\scriptsize, fill=black!8},
    arr/.style={<->, thick, gray!60}
]
\node[layer, minimum width=5cm] (agg) at (0,2.8) {\textbf{Агрегиране}\\Интерпр.\ $\cdot$ Риск};
\node[layer, minimum width=2.2cm] (claim) at (-1.55,1.4) {\textbf{Твърдения}\\Съставяне};
\node[layer, minimum width=2.2cm] (verif) at (1.55,1.4) {\textbf{Проверка}\\Избор};
\node[layer, minimum width=5cm] (did) at (0,0) {\textbf{DID слой}\\Идентичност $\cdot$ Правила $\cdot$ Оценки};

\draw[arr] (agg.south west) -- (claim.north);
\draw[arr] (agg.south east) -- (verif.north);
\draw[arr] (claim.east) -- (verif.west);
\draw[arr] (claim.south) -- (did.north west);
\draw[arr] (verif.south) -- (did.north east);
\end{tikzpicture}
\caption{Четирислойна архитектура на RSN.}
\label{fig:architecture}
\end{figure}

\subsection{Роли на участниците}

Една проверителна транзакция включва до шест различни DID роли (Фиг.~\ref{fig:roles}): \textbf{Издател} ($DID_I$, създава и подава твърдението), \textbf{Субект} ($DID_S$, DID, за когото е твърдението---може да съвпада с Издателя), \textbf{Авторитет} ($DID_A$, подкрепя твърдението срещу такса; може също да предлага свидетелски услуги---\emph{по избор}), \textbf{Свидетел} ($DID_{W_{1..n}}$, инкогнито наблюдател на честността на проверителя чрез challenge кодове---\emph{по избор}), \textbf{Проверител} ($DID_V$, алгоритмично избран да публикува твърдението) и \textbf{RSN} (мрежата, където се публикуват проверените твърдения). Всяка роля може да бъде делегирана на друг DID (Раздел~7.4). Един Авторитет обикновено обединява подкрепянето със свидетелски услуги, но двете функции са логически независими.

\begin{figure}[ht]
\centering
\begin{tikzpicture}[
    role/.style={draw, rounded corners, minimum width=1.1cm, minimum height=0.45cm, align=center, font=\scriptsize, fill=black!8},
    opt/.style={draw, rounded corners, minimum width=1.1cm, minimum height=0.45cm, align=center, font=\scriptsize, fill=black!8, dashed},
    arr/.style={-{Stealth[length=3pt]}, thick},
    oarr/.style={-{Stealth[length=3pt]}, thick, dashed, gray},
    lbl/.style={font=\tiny, fill=white, inner sep=1pt}
]
\node[role] (I) at (0,2.4) {\textbf{Издател}};
\node[role] (S) at (-2.5,2.4) {\textbf{Субект}};
\node[opt] (A) at (2.5,1.2) {\textbf{Авторитет}};
\node[opt] (W) at (-2.5,0) {\textbf{Свидетел}};
\node[role] (V) at (2.5,-0.6) {\textbf{Проверител}};
\node[role, fill=black!5, draw=gray] (N) at (0,-1.8) {RSN};

\draw[arr] (I) -- node[lbl] {за} (S);
\draw[oarr] (I) -- node[lbl,sloped] {подкрепя} (A);
\draw[oarr] (I) -- node[lbl,sloped] {challenge} (W);
\draw[arr] (I) -- node[lbl,sloped] {заявка} (V);
\draw[oarr] (W) -- node[lbl] {наблюдава} (V);
\draw[arr] (V) -- node[lbl] {публикува} (N);
\end{tikzpicture}
\caption{DID роли в проверителна транзакция. Пунктир = по избор (Авторитет, Свидетел).}
\label{fig:roles}
\end{figure}

\subsection{Неподкупност: четири сили}

Неподкупността на RSN не произтича от един-единствен механизъм, а от четири едновременни сили. Нито една не е достатъчна сама по себе си; само тяхната комбинация прави опита за корупция икономически ирационален.

\textbf{Непредсказуема цел.} Проверителят за всяко твърдение се избира недетерминистично чрез consistent hashing (Раздел~5.3). Нападателят не може предварително да набележи конкретен проверител за подкуп, защото самоличността на проверителя е функция на съдържанието на твърдението и предишните итерации, а не на избора на издателя.

\textbf{Репутационен лост---бавно нагоре, бързо надолу.} Репутацията може да се печели само постепенно, чрез трайно последователно поведение. Може обаче да се загуби катастрофално с едно измамно подкрепяне (Раздел~6.2). Тази асиметрия означава, че години натрупана репутация могат да бъдат унищожени с едно нечестно подкрепяне; оптималната стратегия остава да се отхвърлят съмнителните твърдения.

\textbf{Публично обещание.} Всеки притежател на DID публично декларира своята политика---машинно четим набор от правила, които се ангажира да следва. Отклонението между декларацията и действителното поведение е откриваемо и се остойностява като репутационна простъпка, по-тежка от самото нарушение (Раздел~7.3). Следователно лицемерието е по-скъпо от директната нечестност.

\textbf{Асиметрия на цената на атака в мрежата (mesh).} Цената за цензуриране или дестабилизиране на мрежата нараства нелинейно с размера и плътността на мрежата, докато цената за търпимостта към нея остава постоянна. За да се изплати цензурата, централизиран участник би трябвало да я приложи върху цялата мрежа---изисквайки мерки, несъразмерни с всяка мислима полза (Раздел~10).

% ============================================================
% 4. DECENTRALIZED IDENTITY MODEL
% ============================================================
\section{Модел на децентрализирана идентичност}

\subsection{DID със специални свойства}

RSN разширява спецификацията W3C DID Core~\cite{did2022} с: \textbf{Декларирани правила} (машинно четими поведенчески ангажименти), \textbf{Репутационни метаданни} (агрегирана поведенческа оценка) и \textbf{Мрежово маршрутизиране} (променлив onion шлюз, отделен от стабилната позиция в пръстена).

\subsection{Жизнен цикъл на твърдението}

Едно твърдение преминава през шест етапа (Фиг.~\ref{fig:lifecycle}): съставяне, незадължително подкрепяне от авторитет, избор на проверител чрез consistent hashing, решение за проверка (приемане или отхвърляне с итерация), публикуване и репутационно обновяване за всички страни.

\begin{figure}[ht]
\centering
\begin{tikzpicture}[
    stp/.style={draw, rounded corners=2pt, minimum width=1.3cm, minimum height=0.45cm, align=center, font=\tiny, fill=black!5},
    arr/.style={-{Stealth[length=3pt]}, thick},
    lbl/.style={font=\tiny, fill=white, inner sep=1pt}
]
\node[stp] (comp) at (0,0) {Съставяне\\твърдение};
\node[stp, dashed] (auth) at (2.0,0) {Подкрепа\\авторитет};
\node[stp] (hash) at (4.0,0) {Хеш и\\пръстен};
\node[stp] (ver) at (2.0,-1.6) {Проверка\\от провер.};
\node[stp, double] (pub) at (0,-1.6) {Публик.};
\node[stp] (rej) at (4.0,-1.6) {Отхвърл.\\цена $\uparrow$};

\draw[arr] (comp) -- (auth);
\draw[arr] (auth) -- (hash);
\draw[arr] (hash) -- (ver);
\draw[arr] (ver) -- node[lbl] {\checkmark} (pub);
\draw[arr] (ver) -- node[lbl] {$\times$} (rej);
\draw[arr] (rej) -- ++(0,0.8) -| (hash);
\end{tikzpicture}
\caption{Жизнен цикъл на твърдението с итерационен цикъл.}
\label{fig:lifecycle}
\end{figure}

\subsection{Връзка с W3C DID Core}

Моделът на идентичност на RSN е същинско надмножество на спецификацията W3C DID Core~\cite{did2022}. Всички RSN DID са валидни W3C DID. Специфичните за RSN разширения се кодират като свойства на DID документа, дефинирани в отделна спецификация на DID метод, съвместима със съществуваща инфраструктура като ION~\cite{buchner2021} и Ceramic~\cite{ceramic2023}.

% ============================================================
% 5. VERIFICATION AND PROPAGATION
% ============================================================
\section{Проверка и разпространение на информация}

\subsection{Цена на въвеждане на информация}

Основният икономически принцип на RSN е асиметрията между четенето и записа. Четенето на репутационни данни клони към нулева пределна цена за участниците, които са част от DID мрежата и имат достъп, съизмерим с репутацията им---новодошлите трябва постепенно да заслужат по-широк достъп (вж.\ anti-leeching). Записът---разбиран като трайно материализиране на репутацията в мрежата, а не като еднократен технически акт---изисква проверяема кумулативна инвестиция в три измерения: \textbf{време} (години от живота, прекарани в последователно поведение, изпълнени ангажименти и изградени взаимоотношения), \textbf{енергия} (усилие, вложено в честни отношения, грижа за партньорствата и дългосрочна благонадеждност) и \textbf{пари} (инвестиция в собственото име---професионално развитие, качеството на работата, обезщетение за причинени вреди). Репутацията не може да се купи мигновено---тя е резултат от натрупване през целия живот, което системата просто прави видимо и прехвърляемо между контексти. Еднократните технически разходи на протокола (криптографски подписи, такси на проверителите) са пренебрежими в сравнение с тази основополагаща инвестиция.

\subsection{Социален граф и дълбочина на контактите}

RSN е по същество социална мрежа: всеки DID добавя контакти (други DID), които разрешават връзката. Тези контакти имат свои собствени контакти и така нататък. Алгоритмичното търсене на проверители работи в рамките на конфигурируема дълбочина $h$ на свързаните контакти (напр.\ $h=3$: преките контакти, техните контакти и контактите на контактите). Тази архитектура не изисква единен глобален блокчейн---тя естествено благоприятства възникването на общности/клъстери с припокривания в други общности. Всеки участник с DID трябва да има публикувана \textbf{политика}---машинно четим набор от правила, управляващи как се оценяват входящите заявки. Нарушенията на политиката се записват в мрежата като негативни артефакти.

\subsection{Алгоритмичен избор на проверители}

Протоколът за проверка използва consistent hashing~\cite{karger1997} (Фиг.~\ref{fig:verifier}). Проверката е платена услуга: проверителите работят срещу такса, създавайки пазар, в който конкуренцията определя ценообразуването, скоростта и надеждността.

\textbf{Стъпка 1.} Документът на твърдението се конкатенира с хеш резултата от предишната итерация (или $\varnothing$ при първата итерация) и се хешира:
\begin{equation}
\text{pos} = f_h(\text{claim} \| \text{prev\_hash})
\label{eq:hash}
\end{equation}

Алгоритъмът трябва да включва \emph{пълния път} на всички предишни заявки и отговори---не просто хеш на предишната итерация. Пълният отговор на всеки проверител (приемане, отхвърляне, обявена цена, причина) се добавя към нарастващия документ, проверяем чрез криптографски подписи на всяка стъпка.

\textbf{Стъпка 2.} Хешът се съпоставя с $N$ позиции върху пръстена на consistent hashing, равномерно разпределени (напр.\ за $N=4$: $+0^{\circ}, +90^{\circ}, +180^{\circ}, +270^{\circ}$). От всяка позиция търсене по посока на часовниковата стрелка избира най-близкия DID като кандидат-проверител (Фиг.~\ref{fig:multipoint}). $N$ е променлив параметър, който може да зависи от процента на текущо активните DID, часа от денонощието или историческата отзивчивост на мрежата.

\textbf{Стъпка 3.} Проверителят приема (публикува, залагайки репутация) или отхвърля (връщайки се към Стъпка~1 с хеша на отхвърлянето). Когато възел не отговори въпреки декларирания онлайн статус, следващата заявка трябва да включва всички отговори от всичките $N$ предишни кандидат-проверители.

\textbf{Anti-leeching.} Целевите DID могат да поставят такси или ограничения за достъп при заявки от нови DID с малка репутация. Този градуиран (не двоичен) механизъм принуждава новодошлите да изградят репутация чрез истинска активност, преди да потребяват свободно данните на другите.

\begin{figure}[ht]
\centering
\begin{tikzpicture}[
    box/.style={draw, rounded corners=2pt, minimum width=1.3cm, minimum height=0.45cm, align=center, font=\scriptsize, fill=black!5},
    dec/.style={draw, diamond, aspect=2.2, minimum width=0.9cm, align=center, font=\scriptsize, fill=black!5},
    arr/.style={-{Stealth[length=3pt]}, thick},
    lbl/.style={font=\tiny, fill=white, inner sep=1pt}
]
\node[box] (iss) at (0,0) {Издател};
\node[box] (hash) at (0,-1.1) {$f_h$(claim $\|$\\prev\_hash)};
\node[box] (ring) at (0,-2.2) {Пръстен\\по часовн.};
\node[box, dashed] (spam) at (0,-3.2) {Анти-спам\\депозит};
\node[dec] (check) at (0,-4.4) {Проверка\\политика};
\node[box] (rej) at (2,-5.6) {Следваща\\итерация};
\node[box] (fee) at (0,-5.6) {Крайна такса =\\$f$(данни, реп., пол.)};
\node[box, double] (pub) at (0,-6.8) {Публикувано};

\draw[arr] (iss) -- (hash);
\draw[arr] (hash) -- (ring);
\draw[arr] (ring) -- (spam);
\draw[arr] (spam) -- (check);
\draw[arr] (check) -- node[lbl] {$\times$} (rej);
\draw[arr] (check) -- node[lbl] {\checkmark} (fee);
\draw[arr] (fee) -- (pub);
\draw[arr] (rej.east) -- ++(0.5,0) |- (hash.east);
\end{tikzpicture}
\caption{Протокол за избор на проверители. Анти-спам депозитът е по избор и може да бъде възстановим. Крайната такса се плаща само при приемане.}
\label{fig:verifier}
\end{figure}

Всяко отхвърляне добавя пълния отговор на проверителя (включително причини и условия) към документа на твърдението, разширявайки съдържанието му. От разширения документ се изчислява недетерминистично нов хеш, който се съпоставя с различна позиция в пръстена и избира различен проверител. Документирането на пълния път е задължително---издателят не може да предскаже или повлияе на следващия проверител. Ако проверител пренебрегне заявка въпреки съвместима декларирана политика, свидетелите (ако присъстват) записват това и издателят може да приложи свидетелски доказателства при окончателното публикуване.

\begin{figure}[ht]
\centering
\begin{tikzpicture}[scale=0.65]
% Ring
\draw[thick] (0,0) circle (2.2cm);
% DID nodes on ring
\foreach \a in {10,55,80,120,150,195,230,260,305,340} {
  \fill[black!40] (\a:2.2) circle (1.5pt);
}
% Hash offset arrow pointing to initial position
\draw[-{Stealth[length=3pt]}, thick, black!60] (0,0) -- node[font=\tiny,above,sloped,pos=0.6] {$f_h$} (37:2.2);
\fill[black] (37:2.2) circle (2pt);
\node[font=\tiny,anchor=south west] at (37:2.35) {$h_0$};
% N=4 points offset from h0 by +0, +90, +180, +270
\foreach \offset/\l in {0/P1,90/P2,180/P3,270/P4} {
  \pgfmathsetmacro{\ang}{37+\offset}
  \node[fill=black, circle, inner sep=1.5pt] (p\offset) at (\ang:2.2) {};
  \node[font=\tiny,font=\bfseries] at (\ang:2.75) {\l};
  % clockwise search arc
  \draw[-{Stealth[length=2.5pt]}, thick, gray] (\ang:2.2) arc (\ang:\ang+30:2.2);
}
\node[font=\scriptsize, align=center] at (0,0) {Хеш\\пръстен};
\node[font=\tiny, align=center] at (0,-3.2) {$h_0 = f_h(\text{claim})$, след това $+0^{\circ},+90^{\circ},+180^{\circ},+270^{\circ}$\\Всяка точка $\rightarrow$ търсене по часовниковата стрелка};
\end{tikzpicture}
\caption{Многоточков избор ($N{=}4$). Хешът $h_0$ задава отместването; 4 точки равномерно разпределени от $h_0$.}
\label{fig:multipoint}
\end{figure}

\subsection{Свидетелски протокол}

Ролята \textbf{Свидетел} осигурява инкогнито отчетност за честността на проверителя чрез криптографски протокол с challenge код:

\textbf{Стъпка W1.} Заявителят подписва заявката за проверка и я изпраща на $N_w$ свидетели---контакти от социалната мрежа на заявителя или специализирани доставчици на свидетелски услуги, работещи срещу такса.

\textbf{Стъпка W2.} Всеки свидетел $w_i$ подписва цялата подписана заявка, запазва оригиналния подпис $\sigma_i$ и изчислява challenge код $c_i = h(\sigma_i)$. Challenge кодът се добавя към заявката.

\textbf{Стъпка W3.} Заявката, вече носеща $N_w$ challenge кода, се изпраща на проверителя. Проверителят вижда кодовете, но \emph{не може да определи} кои са свидетелите, нито дали кодовете съответстват на реални подписи.

\textbf{Стъпка W4 (честен проверител).} Ако проверителят отговори според декларираната си политика, challenge кодовете остават непрозрачни. Свидетелите никога не се разкриват. Не се публикуват допълнителни данни.

\textbf{Стъпка W5 (нарушение на политиката).} Ако проверителят наруши политиката си (пренебрегне заявката, отговори непоследователно), заявителят разполага с оригиналните подписи на свидетелите $\sigma_i$. Те могат да бъдат публикувани като придружаващи данни: всеки може да провери $c_i = h(\sigma_i)$, доказвайки, че свидетелят е присъствал. Заявителят може да публикува независимо---проверителят вече е подписал challenge кодовете в отговора си.

\textbf{Свойството блъф.} Заявителят може да замести някои или всички challenge кодове със случайни стойности. Проверителят не може да различи истинските кодове (подкрепени от авторитети с висока репутация) от шум. Тази \emph{несигурност} е механизмът за принудително изпълнение: всяка заявка носи риска, че уважаван авторитет наблюдава инкогнито. Цената за поддържане на този натиск е близка до нулата (генериране на случайни числа), докато потенциалната цена на нечестността за проверителя е катастрофална. Честното поведение е стимулирано дори при липсата на реални свидетели.

\textbf{Авторитет като свидетел.} Авторитет, който подкрепя дадено твърдение, може да обедини свидетелски услуги: „Подкрепям твоето твърдение и служа като инкогнито свидетел по време на проверката.“ Това е естествен бизнес модел---Авторитетът вече разполага с контекст. При все това двете роли са логически независими.

\subsection{Принцип на скъпия радикализъм}

Три разходни канала се усилват едновременно (Фиг.~\ref{fig:radicalism}):

\textbf{Канал 1: Цена на итерация.} Всяко отхвърляне налага нова итерация, консумираща време и ресурси. Линеен растеж.

\textbf{Канал 2: Разрастване на документа.} Всяка итерация добавя пълния отговор на проверителя към документа на твърдението. Растежът е линеен, но с базово отместване: първоначалният полезен товар (съдържание на твърдението, подкрепа от авторитет, криптографски подписи) вече има нетривиален размер $B_0$. Общ размер след $k$ итерации: $S(k) = B_0 + k \cdot \Delta$, където $\Delta$ е средният размер на отговора.

\textbf{Канал 3: Рискова премия на проверителя.} Рискът е субективен. Проверителят преценява дали декларираните политики на заявяващия DID влизат в конфликт със собствените търговски, социални или политически интереси на проверителя. Премията може да ескалира експоненциално с възприеманото разстояние от „идеала“ на проверителя: $P(d) = \alpha \cdot e^{\beta d}$, където $d$ е субективната метрика на разстояние на проверителя. Отвъд лична граница на непреминаване проверителят може да откаже изцяло.

\begin{figure}[ht]
\centering
\begin{tikzpicture}
\begin{axis}[
    width=\columnwidth,
    height=4.5cm,
    xlabel={\footnotesize Радикализъм},
    ylabel={\footnotesize Относителна цена (log)},
    ymode=log,
    xmin=0, xmax=100,
    ymin=1, ymax=10000,
    grid=major,
    grid style={gray!20},
    legend style={at={(0.03,0.97)},anchor=north west,font=\tiny,draw=gray!50},
    tick label style={font=\tiny},
    label style={font=\footnotesize},
]
\addplot[black, thick, domain=0:100, samples=50] {1 + 0.3*x};
\addlegendentry{Цена на итерация (линейна)}
\addplot[black, thick, dashed, domain=0:100, samples=50] {1 + 0.15*x};
\addlegendentry{Разрастване на документа (линейно)}
\addplot[black, thick, dotted, line width=1.2pt, domain=0:100, samples=100] {exp(0.08*x)};
\addlegendentry{Рискова премия (експоненциална)}
\end{axis}
\end{tikzpicture}
\caption{Ескалация на цената през три канала. Рисковата премия доминира при висок радикализъм.}
\label{fig:radicalism}
\end{figure}

Критично важно е, че това не е цензура. Нито едно твърдение не е забранено. Системата остойностява риска (Табл.~\ref{tab:costs}).

\begin{table}[ht]
\centering
\caption{Сравнение на разходите за различни типове твърдения.}
\label{tab:costs}
\footnotesize
\begin{tabular}{@{}lcccc@{}}
\toprule
\textbf{Тип} & \textbf{Итер.} & \textbf{Док.} & \textbf{Такса} & \textbf{Общо} \\
\midrule
Достоверно & $k_{\min}$ & $B_0$ & $P_{\min}$ & Ниско \\
Спорно & $k_{\text{med}}$ & $B_0{+}k\Delta$ & $\alpha e^{\beta d}$ & Умерено \\
Радикално & $k_{\max}$ & $\gg B_0$ & $\gg P_{\min}$ & Много високо \\
Измамно & $\to\infty$ & --- & --- & Непубликуемо \\
\bottomrule
\end{tabular}
\end{table}

% ============================================================
% 6. REPUTATION AND RISK
% ============================================================
\section{Оценка на репутация и риск}

\subsection{Модел на натрупване на репутация}

Репутацията проявява три свойства: \textbf{бавно натрупване} (постепенен растеж чрез последователно поведение), \textbf{бързо разрушаване} (една-единствена измама може катастрофално да понижи оценката) и \textbf{индивидуална оценка} (всяка потребяваща страна може да прилага собствена агрегираща функция).

\subsection{Репутационни авторитети}

Един репутационен авторитет преглежда твърденията и, ако е удовлетворен, ги подписва съвместно---залагайки собствената си репутация (Фиг.~\ref{fig:authority}). Асиметрията е по замисъл: печеленето на репутация е бавно (постепенно $+\delta$ на честно подкрепяне), докато загубата ѝ е катастрофална ($-\Delta \gg \delta$ на фалшиво подкрепяне; илюстративно, напр.\ $+2\%$ срещу $-70\%$). Оптималната стратегия винаги е да се отхвърлят съмнителните твърдения. Конкретните стойности на $\delta$ и $\Delta$ са системни параметри.

\begin{figure}[ht]
\centering
\begin{tikzpicture}[
    box/.style={draw, rounded corners=2pt, minimum width=2cm, minimum height=0.6cm, align=center, font=\scriptsize, fill=black!5},
    arr/.style={-{Stealth[length=4pt]}, thick}
]
\node[box] (exam) at (0,0) {Авторитет проверява\\Репутация: $R$};
\node[box] (true) at (-2,-1.3) {Истина: $R{\to}R{+}\delta$\\Повече клиенти};
\node[box] (false) at (2,-1.3) {Лъжа: $R{\to}R{-}\Delta$\\Клиентите си тръгват};
\node[box] (insight) at (0,-2.6) {Печалба: бавна ($+\delta$)\\Загуба: мигновена ($-\Delta$)};

\draw[arr] (exam) -- node[left,font=\tiny] {истина} (true);
\draw[arr] (exam) -- node[right,font=\tiny] {лъжа} (false);
\draw[arr] (true) -- (insight);
\draw[arr] (false) -- (insight);
\end{tikzpicture}
\caption{Репутационен авторитет---асиметрични стимули.}
\label{fig:authority}
\end{figure}

\subsection{Многомерна репутация}

Репутацията не е скалар, а вектор в множество измерения: финансова благонадеждност, лична почтеност, професионална компетентност, гражданска ангажираност и други (Фиг.~\ref{fig:reputation-radar}). Различните доставчици на оценка проектират този вектор различно в зависимост от контекста---заемодател приоритизира финансовата благонадеждност, работодател професионалната компетентност, съсед гражданското поведение. Няма единна „правилна“ репутационна оценка; само гледни точки.

\begin{figure}[ht]
\centering
\begin{tikzpicture}[scale=0.55]
\foreach \a/\l in {90/Финансова,162/Почтеност,234/Профес.,306/Гражданска,18/Социална} {
  \draw[gray!40] (0,0) -- (\a:2.5);
  \node[font=\tiny, align=center] at (\a:3.1) {\l};
  \foreach \r in {0.8,1.6,2.4} {
    \fill[black!15] (\a:\r) circle (0.5pt);
  }
}
\draw[thick, fill=black!10] (90:2.0) -- (162:1.2) -- (234:1.8) -- (306:0.9) -- (18:1.5) -- cycle;
\draw[thick, dashed] (90:1.4) -- (162:2.1) -- (234:0.8) -- (306:2.0) -- (18:1.1) -- cycle;
\node[font=\tiny] at (0,-3.3) {Плътно: DID A \quad Пунктир: DID B};
\end{tikzpicture}
\caption{Многомерни репутационни вектори. Различните DID проявяват различни профили в различните измерения.}
\label{fig:reputation-radar}
\end{figure}

\subsection{Демократизирана оценка на риска}

RSN демократизира систематичната оценка на риска---способност, понастоящем концентрирана във финансовите институции, но приложима далеч отвъд банковото дело. Промишлени конгломерати, оценяващи надеждността на доставчици, застрахователни компании, остойностяващи поведенчески риск, надлежна проверка при сливания и придобивания, оценка на срока на живот на оборудване в производството---навсякъде, където рискът на насрещната страна изисква оценка, RSN предлага децентрализирана, пазарно управлявана алтернатива. Пазар от интерпретационни услуги превежда суровите многомерни репутационни данни в приложими обобщения, правейки оценката на насрещната страна достъпна за всеки участник на пределна цена.

% ============================================================
% 7. CONSENSUS AND DISPUTE RESOLUTION
% ============================================================
\section{Консенсус и разрешаване на спорове}

\subsection{Модел на децентрализирана справедливост}

RSN преформулира справедливостта като проблем на двустранно договаряне. Заплахата от постоянна, проверяема репутационна вреда е по-ефективен възпиращ фактор от съдебно производство, което ощетената страна не може да си позволи.

\subsection{Емергентен консенсус}

Всеки участник поддържа личен праг на удовлетворение. Съвкупният консенсус на мрежата възниква като статистическа средна стойност на хиляди двустранни договаряния (Фиг.~\ref{fig:consensus}). Отговор далеч над консенсуса (варварски) е скъп за публикуване. Отговор близо до консенсуса (пропорционален) е евтин. Консенсусът не е правило---той е ценови сигнал.

\begin{figure}[ht]
\centering
\begin{tikzpicture}[
    person/.style={draw, rounded corners=2pt, minimum width=1.5cm, minimum height=0.5cm, align=center, font=\scriptsize, fill=black!5},
    arr/.style={-{Stealth[length=4pt]}, thick}
]
\node[person] (a) at (-2,0) {Строг\\(висок)};
\node[person] (b) at (0,0) {Прагматичен\\(среден)};
\node[person] (c) at (2,0) {Прощаващ\\(нисък)};
\node[person, fill=black!10, minimum width=3cm] (cons) at (0,-1.2) {Мрежов консенсус\\= стат.\ средна};
\node[person] (exp) at (-2,-2.4) {Варварски\\скъп};
\node[person, double] (prop) at (0,-2.4) {Пропорц.\\евтин};
\node[person] (neg) at (2,-2.4) {Небрежен\\скъп};

\draw[arr] (a) -- (cons);
\draw[arr] (b) -- (cons);
\draw[arr] (c) -- (cons);
\draw[arr] (cons) -- (exp);
\draw[arr] (cons) -- (prop);
\draw[arr] (cons) -- (neg);
\end{tikzpicture}
\caption{Емергентен консенсус от индивидуалните прагове на удовлетворение.}
\label{fig:consensus}
\end{figure}

\subsection{Калибриране на наказанието}

Всеки притежател на DID публично декларира отговори на конкретни категории простъпки. Непропорционално строгите или снизходителните декларации привличат негативни репутационни сигнали. Пропорционалните декларации се приемат без триене---самокалибрираща се наказателна система.

Консенсусните декларации сами по себе си подлежат на откриване на лицемерие: декларативната ангажираност с консенсусна позиция при непоследователно поведение представлява репутационна простъпка, по-тежка от самото отклонение, тъй като демонстрира умишлена измама.

\subsection{Делегиране}

Всички дейности в рамките на един DID---с едно изключение---могат да бъдат делегирани на друг DID. \textbf{Ролята Субект---DID, за чието поведение е твърдението---не може да бъде делегирана; тя е непрехвърляемо обвързана с притежателя на идентичността, за когото се отнася твърдението.} Другите активни роли---Издател, Авторитет, Свидетел, Проверител---могат да бъдат прехвърлени на определен делегиран DID, било за проверка, подаване на твърдения, участие в консенсус или разрешаване на спорове. Делегирането е отменимо по всяко време. Това дава възможност за специализация (напр.\ професионални проверителни услуги), докато притежателят запазва крайния контрол и отговорност. Делегирането се прилага в цялата система и е съществено за мащабируемостта.

\subsection{От индивидуален морал към емергентен обществен договор}

Декларираната политика на всеки DID представлява формализация на индивидуалния морал: какво притежателят смята за приемливо, как отговаря на конкретни поведения, какво изисква от насрещните страни. Тези политики не се налагат от проектант на системата---те се избират от всеки участник и се изразяват в машинно четима форма.

Тази формализация дава възможност за тристепенно възникване. Първо, индивидуалният морал се кодира като \textbf{политика}---набор от декларирани правила, прагове и функции на отговор, които DID се ангажира да следва. Второ, съвкупността от всички политики в мрежата произвежда \textbf{емергентна етика}---статистически наблюдаеми норми, които никой отделен участник не е проектирал, но които възникват от взаимодействието на хиляди индивидуални морални позиции~\cite{durkheim1893}. Трето, тази емергентна етика, изразена като споделени поведенчески норми, налагани чрез двустранни ценови сигнали, съставлява \textbf{измерим обществен договор}---не теоретична конструкция, която никой не е подписвал~\cite{rawls1971}, а емпирично наблюдаем набор от правила, подкрепен от действително поведение.

Този процес отразява находки от теорията на моралните основи~\cite{haidt2012}: човешкият морал е интуитивен, плуралистичен и културно променлив. RSN не изисква морален консенсус като вход; тя произвежда поведенчески консенсус като изход. Получаващият се обществен договор не е статичен---той се развива, докато участниците се присъединяват, напускат и коригират политиките си в отговор на мрежовата обратна връзка. За разлика от класическата теория на обществения договор, този договор е отменим, наблюдаем и приложим без суверен.

% ============================================================
% 8. ECONOMIC MODEL
% ============================================================
\section{Икономически модел}

Предходните раздели описаха инструмент---проверителните механизми на RSN, структурата на разходите и емергентния консенсус. Този раздел описва методология за прилагане на този инструмент към фискален и управленски преход. Инструментът е с общо предназначение; методологията по-долу е едно от възможните приложения.

\subsection{Структура на стимулите}

Репутацията като капитал създава преки стимули за честно поведение. При нормална работа участниците изграждат репутация естествено чрез ежедневни транзакции и изпълнени ангажименти. Основният разход е за \emph{негативни} твърдения---записване на вреда, причинена от други. Тази асиметрия стимулира полезно, надеждно поведение: да бъдеш честен не струва нищо допълнително, докато да бъдеш вреден създава документирана следа, която другите ще платят, за да съхранят. Лицемерието---деклариране на правила без да ги следваш---е най-скъпият режим на провал, тъй като демонстрира умишлена измама.

\subsection{Паралелна фискална система}

Три компонента дават възможност за конкурентен натиск: (1)~\textbf{Прост данък}---единна пропорционална ставка без изключения, първоначално незначително под съществуващата ефективна ставка; (2)~\textbf{Електронна регистрация на разходите (ESR)}---обръщайки чешкия модел EET, така че гражданите да наблюдават държавните разходи; (3)~\textbf{Гражданско направляемо разпределение на данъци}---започвайки от няколко процента през първата година и достигайки след десетилетие някъде от ниските двуцифрени числа до пълна миграция към система, направлявана от гражданите, в зависимост от темпа на приемане. Действителното темпо зависи от скоростта, с която възникват свободнопазарни алтернативи на държавните услуги, и от готовността на държавата да сътрудничи на прехода; функцията на времето не е задължително линейна и няма причина да се поставя горна граница на това докъде в крайна сметка може да достигне приемането.

% ============================================================
% 9. MIGRATION PATH
% ============================================================
\section{Път на миграция}

Миграцията протича през три фази (Фиг.~\ref{fig:migration}): \textbf{Фаза~1: Наслагване} (RSN като информационен слой, държавата е единствен доставчик), \textbf{Фаза~2: Конкуренция} (RSN предоставя функционални алтернативи, използване на двойна система) и \textbf{Фаза~3: Преход} (RSN превъзхожда държавните еквиваленти, доброволна миграция). Преходът е обратим на всеки етап.

\begin{figure}[ht]
\centering
\begin{tikzpicture}[
    phase/.style={draw, rounded corners=3pt, minimum width=1.8cm, minimum height=0.8cm, align=center, font=\scriptsize, fill=black!5},
    arr/.style={-{Stealth[length=5pt]}, thick}
]
\node[phase] (p1) at (0,0) {\textbf{Фаза 1}\\Наслагване};
\node[phase] (p2) at (2.2,0) {\textbf{Фаза 2}\\Конкуренция};
\node[phase] (p3) at (4.4,0) {\textbf{Фаза 3}\\Преход};

\draw[arr] (p1) -- (p2);
\draw[arr] (p2) -- (p3);
\draw[arr, dashed, gray] (p3.south) -- ++(0,-0.6) -| node[below,font=\tiny,pos=0.25] {връщане} (p2.south);
\end{tikzpicture}
\caption{Тристепенна миграция---обратима на всеки етап.}
\label{fig:migration}
\end{figure}

Основният механизъm на натиск е фискален: когато условните данъкоплатци достигнат „икономически значимо малцинство $X$“---група, достатъчно голяма, че репресиите срещу нея да причинят нетривиални икономически щети и гражданското неподчинение да се превърне в достоверна заплаха---държавата влиза в преговори. За илюстрация използваме $X \approx 15\%$ от БВП, но действителният праг зависи от конкретната икономика.

% ============================================================
% 10. SECURITY ANALYSIS
% ============================================================
\section{Анализ на сигурността}

Разглеждаме пет категории противници: индивидуални недобросъвестни участници, организирани групи, корпоративни противници, противници на държавно ниво и противници на ниво протокол.

\textbf{Устойчивост срещу Sybil}~\cite{douceur2002}. Изграждането на репутация изисква трайно, проверяемо взаимодействие. Цената на успешна Sybil атака нараства линейно с фалшивите идентичности и квадратично с целевото ниво на репутация. Оперирането на множество DID паралелно е възможно, но скъпо по замисъл---всяка идентичност трябва независимо да натрупа репутация чрез истинска активност. В свободните общества тази цена възпира несериозната злоупотреба; в диктатурите паралелните DID се превръщат в инструмент за оцеляване, който дава възможност за компартментализирана съпротива, подземна координация и по-безопасно навигиране в черния пазар.

\textbf{Защита срещу тайно споразумение.} Модели на взаимно подкрепяне сред затворени групи са откриваеми чрез анализ на социалния граф. Функциите за агрегиране на репутацията обезценяват твърдения от тясно клъстеризирани източници.

\textbf{Противник на държавно ниво.} Onion маршрутизиране, отсъствие на централизирана инфраструктура и разпределение на ролята Проверител сред цялата база от участници гарантират, че потискането изисква мерки, несъразмерни с всяка полза.

\textbf{Поверителност срещу отчетност.} Псевдонимността---средно положение между анонимност и разкриване на идентичност---позволява на DID да натрупват значима репутация, без да разкриват реалната самоличност на притежателя.

% ============================================================
% 11. PRIVACY
% ============================================================
\section{Съображения за поверителност}

Моделът на поверителност на RSN е изграден върху псевдонимност. Притежателят контролира разкриването на идентичността: минимално (чист псевдоним), избирателно (свързани конкретни удостоверения) или пълно (асоциирана правна идентичност). Публикуваните твърдения са постоянни---системата поддържа контекстуална корекция, но не изтриване. Притежателят може да изостави компрометиран DID и да започне наново, губейки натрупаната репутация.

% ============================================================
% 12. RELATED WORK
% ============================================================
\section{Свързани работи}

Спецификацията W3C DID Core~\cite{did2022} и мрежата Ceramic~\cite{ceramic2023} осигуряват основата на идентичността. EigenTrust~\cite{kamvar2003} демонстрира разпределено агрегиране на репутация без централен авторитет. SBT~\cite{weyl2022} въведоха непрехвърляеми социални токени. RSN се различава по акцента си върху икономическата цена като филтър за качество и по механизма си за остойностяване на радикализма.

DAO~\cite{buterin2014} и квадратичното гласуване~\cite{buterin2019} демонстрираха осъществимостта на децентрализираното управление. RSN извежда консенсус от двустранни ценови преговори, а не от гласуване.

Предложението се опира на анархо-капиталистическата теория~\cite{rothbard1973,hoppe2001} за частно предоставяне на услуги, но се отклонява в две критични отношения: то проектира за злоба и възмездни импулси, вместо да предполага рационалност, и предлага постепенен преход, а не революция. Концепцията за емергентни обществени договори има предшественици в теорията на Хайек за спонтанния ред~\cite{hayek1973}.

\textbf{Анархо-агоризъм.} RSN споделя значителна обща основа с агористката контра-икономика~\cite{konkin2006}---особено акцента върху изграждането на паралелни институции и доброволна размяна. Агоризмът обаче е склонен да предполага рационалния личен интерес като достатъчен механизъм за координация и често му липсва конкретен път на миграция от съществуващите държавни структури. RSN изрично включва нерационалните поведенчески склонности и осигурява механизъм за фискален натиск за постепенен преход.

\textbf{Меритокрация.} RSN може да представлява първо функционално описание на това как меритокрацията~\cite{young1958} може да възникне като системно свойство, а не като лозунг. Традиционните меритократични системи се провалят, защото изискват централен авторитет, който да дефинира и налага „заслугата“. В RSN заслугата възниква от двустранни оценки: тези, които доставят стойност, натрупват репутация; никакъв комитет не решава кой има заслуга.

\textbf{Собствеността като привилегия.} RSN се отклонява фундаментално от анархо-капиталистическото и австрийско-школовото третиране на правата на собственост като естествени и абсолютни. В рамката на RSN собствеността е \emph{привилегия}---социално признание, което в крайни случаи може да бъде оттеглено от общността, когато участник фундаментално наруши споделените норми (предателство, отказ да защити общността в екзистенциална криза, систематична експлоатация). Това не е държавна експроприация, а оттегляне на социалното признание от мрежата на двустранните взаимоотношения.

% ============================================================
% 13. LIMITATIONS
% ============================================================
\section{Обсъждане и ограничения}

\textbf{Студен старт.} Стойността на RSN зависи от мрежовите ефекти; ранните участници се сблъскват с ограничена стойност, докато участието достигне праг. При все това, дори от ранните етапи, DID мрежата служи като полезен допълнителен източник на информация. Очаква се ранната оценка на репутацията и авторитетите да се развива постепенно с нарастваща изтънченост.

\textbf{Anti-leeching и контрол на достъпа.} Целевите DID могат да налагат градуирани такси и ограничения за достъп при заявки от DID с ниска репутация. Това не е двоично, а непрекъсната скала: колкото по-малко репутация има заявяващият DID, толкова по-малко информация получава, толкова по-дълго чака и толкова повече плаща. Този механизъм стимулира истинското участие пред пасивното потребление.

\textbf{Неравенство в достъпа.} Цената на публикуване на твърдения може да отсее легитимни оплаквания от икономически по-необлагодетелствани участници. Смекчаване: всеки участник едновременно служи като потенциален проверител (или делегира тази роля) и печели проверителни такси. През достатъчен времеви хоризонт нерадикален участник би трябвало да се приближи до финансова неутралност---приходите от проверка приблизително компенсират разходите за публикуване. Това е статистическо очакване, а не гаранция.

\textbf{Репутационно неравенство.} Ранните участници натрупват предимства, които могат да създадат бариери за закъснелите. При все това оценката на репутацията се извършва от доставчици трети страни, които могат да изберат да смекчат предимството на първия ход чрез своите агрегиращи алгоритми. Да бъдеш пръв с добра репутация е легитимно сравнително икономическо предимство---и това е по замисъл.

\textbf{Отворени въпроси.} Оптималните разходни функции, стандартите за агрегиране, управлението на протокола и взаимодействието с генерирани от ИИ синтетични репутационни данни остават области за бъдеща работа.

Настоящата статия не твърди, че RSN ще произведе оптимални резултати при всички обстоятелства. Тя твърди, че RSN представлява \emph{осъществима} алтернатива---технически реализуема, икономически самоиздържаща се и социално съвместима с човешките поведенчески склонности, каквито те действително са.

% ============================================================
% 14. CONCLUSION
% ============================================================
\section{Заключение}

Настоящата статия представи Репутационната социална мрежа---децентрализиран протокол, който дава възможност за псевдонимен, проверяем обмен на репутация. Принципът на скъпия радикализъм гарантира, че измамните твърдения се остойностяват до непубликуемост без цензура. Предложеният път на миграция дава възможност за постепенен, обратим преход от съществуващите институции. Проектът включва човешката природа такава, каквато е---включително дребнавост, злоба и желание за възмездно удовлетворение---вместо да изисква идеологическа трансформация. Резултатът не е утопия, а система, в която справедливостта е достъпна, репутацията е заслужена и цената на простъпката се носи от страната, която я е причинила.

% ============================================================
% REFERENCES
% ============================================================
\bibliography{references}

\end{document}
